Una cooperativa agrícola dispone de tres tipos de cultivos: \textbf{trigo}, \textbf{maíz}  y \textbf{girasol}  
Cada uno requiere una cierta cantidad de agua y superficie, y aporta una \textbf{producción total de biomasa (en toneladas)} que la cooperativa quiere \textbf{maximizar}.

Los datos son los siguientes:

\begin{center}
\begin{tabular}{lccc}
\toprule
\textbf{Cultivo} & \textbf{Agua necesaria (m$^3$/ha)} & \textbf{Biomasa (ton/ha)} \\
\midrule
\textbf{Trigo} ($x_1$) & 300 & 4 \\
\textbf{Maíz} ($x_2$) & 500 & 6 \\
\textbf{Girasol} ($x_3$) & 400 & 5 \\
\bottomrule
\end{tabular}
\end{center}

Cada semana, la cooperativa dispone de \textbf{30\,000 m$^3$ de agua} y \textbf{80 hectáreas de terreno total}.  
Además, por exigencias ecológicas, el cultivo de maíz y girasol juntos debe ocupar \textbf{al menos 40 hectáreas}.

El modelo de Programación Lineal que permite determinar el plan de cultivo que maximiza la producción de biomasa es:

\[
\begin{aligned}
\max \quad & z = 4x_1 + 6x_2 + 5x_3 \\
\text{r. t.} \quad 
& 3x_1 + 5x_2 + 4x_3 \le 300 \quad && \text{(disponibilidad de agua)} \\
& x_1 + x_2 + x_3 \le 80 \quad && \text{(superficie total)} \\
& x_2 + x_3 \ge 40 \quad && \text{(requisito ecológico)} \\
& x_1, x_2, x_3 \ge 0
\end{aligned}
\]

donde $x_1$, $x_2$ y $x_3$ representan, respectivamente, las hectáreas dedicadas a trigo, maíz y girasol,

Se conoce la tabla correspondiente a la solución óptima de dicho problema:

\begin{center}
\begin{tabular}{c|c|ccccccc|}
 & $z$ & $x_1$ & $x_2$ & $x_3$ & $h_1$ & $h_2$ & $h_3$ & $a_3$\\ 
      & $-380$ & $0$ & $0$ & $0$ & $-1$ & $-1$ & $0$ & $0$ \\ 
\hline
$x_1$ & 40  & $1$ & $0$ & $0$ & $0$ & $1$ & $1$ & $-1$\\ 
$x_3$ & 20  & $0$ & $0$ & $1$ & $-1$ & $3$ & $-2$ & $2$\\ 
$x_2$ & 20  & $0$ & $1$ & $0$ & $1$ & $-3$ & $1$ & $-1$\\ 
\hline
\end{tabular}
\end{center}
 
 
\bigskip
Se pide lo siguiente:
\begin{enumerate}[label=\alph*)]
    \item \textbf{Interpretar la solución óptima e indicar cuál es el plan de cultivo óptimo de la cooperativa y del cumplimiento de los requisitos descritos.}

    El plan de cultivo óptimo pasa por:
    \begin{itemize}
        \item Producir 40 hectáreas de trigo, 20 hectáreas de maíz y 20 hectáreas de girasol.
        \item Se haría uso de toda el agua.
        \item Se emplearía toda la superficie.
        \item Se cumpliría con el requisito ecológico por el mínimo imprescindible.
    \end{itemize}
    Esto permite generar una biomasa total de 380 toneladas.

    % Impacto sobre el beneficio y sobre el plan de producción del cambio de B3 = 1
    \item \textbf{Cuál sería el impacto sobre el plan de producción y sobre la producción total de biomasa si se incrementase en una unidad el requisito ecológico relativo al mínimo del total de producción de maíz y girasol.}

    Aumentar una unidad el requisito ecológico es equivalente a $h_3=-1$, es decir, que en la situación actual se estuviera incumpliendo el requisito en una unidad. Por lo tanto, el impacto sería el siguiente:

    \begin{itemize}
        \item La producción total de biomasa no se modificaría porque $V^B_{h_3}=0$.
        \item La superficie de trigo sería 39, una unidad menor porque $p^B_{16}=1$.
        \item La superficie de maíz sería 22, dos unidades mayor porque $p^B_{26}=-2$.
        \item La superficie de girasol sería 19, una unidad menor porque $p^B_{36}=1$.
    \end{itemize}
        

    % Sensibilidad de c_1
    \item \textbf{Cuál es la producción de masa por hectárea de maíz, por debajo de la cual no sería interesante cultivar maíz}.

    \[
    \begin{aligned}
    V^B = c-c^Bp^B = 
    \begin{pmatrix}
    4 & c_2 & 5 & 0 & 0 & 0\\
    \end{pmatrix}
     - \begin{pmatrix}
    4 & 5 & c_2\\
    \end{pmatrix}
     \begin{pmatrix}
    1 & 0 & 0 & 0 & 1 & 1\\
    0 & 0 & 1 & -1 & 3 & -2\\
    0 & 1 & 0 & 1 & -3 & 1\\
    \end{pmatrix}
     = \\ = 
    \begin{pmatrix}
    4 & c_2 & 5 & 0 & 0 & 0\\
    \end{pmatrix}
    -
    \begin{pmatrix}
    4 & c_2 & 5 & c_2-5 & 19-3c_2 & c_2-2\\
    \end{pmatrix} = \\
    = \begin{pmatrix}
    0 & 0 & 0 & 5 - c_2 & 3c_2 - 19 & 2 - c_2\\
    \end{pmatrix}
    \end{aligned}
    \]

    La solución actual sigue siendo óptima si $V^B \leq 0$ y, por lo tanto:

    \[
    \begin{aligned}
    5 - c_2 &\le 0, \\
    3c_2 - 19 &\le 0, \\
    2 - c_2 &\le 0.
    \end{aligned}
    \]
    
    Por tanto, la solución actual sigue siendo óptima mientras se cumpla simultáneamente:
    
    \[
    5 \le c_2 \le \tfrac{19}{3}.
    \]
    
    Si la producción de biomasa por hectárea que ofrece el maíz es inferior a 5 la variable $h_1$ tendrá criterio del simplex positivo, y la solución dejará se ser óptima. La nueva solución óptima pasa porque $h_1$ sea básica y $x_2$ deje de serlo (no se cultivaría maíz).

    % Nuevo cultivo
    \item \textbf{Existe la posibilidad de introducir un cuarto cultivo, \textbf{soja}, que genera biomasa.
    Se sabe que requiere \(\,350\,\text{m}^3/\text{ha}\) de agua y produce 4.5 ton/ha.
    La soja no forma parte del requisito ecológico (esto es, no contribuye al mínimo conjunto de maíz y girasol).}

    Se trata de una nueva actividad, para la cual habría que definir una nueva variable $x_4$ con:
    
    \begin{itemize}
    \item $c_{x_4} = 4.5$.
    \item $A_{x_4} =
    \begin{pmatrix}
        3.5 \\[2pt]
        1 \\[2pt]
        0
    \end{pmatrix}$.
    \end{itemize}

    Sería interesante cultivar soja si $V^B_{x_4}>0$:

    \[
    V^B_{x_4} = c_{x_4} - \pi^BA_{x_4} = 4.5 - 
    \begin{pmatrix}
    1 & 1 & 0    
    \end{pmatrix}
    \begin{pmatrix}
        3.5 \\[2pt]
        1 \\[2pt]
        0
    \end{pmatrix} = 4.5-4.5=0
    \]
    Es decir, se podría cultivar soja pero no tendría impacto sobre el total de biomasa generada.

    % Restricción de mercado de un producto (Lemke)
    \item \textbf{Debido a las previsiones meteorológicas para la temporada, la cooperativa ha decidido producir un máximo de 10 hectáreas de maíz, ¿cuál sería el nuevo plan de cultivo óptimo en este caso?}

    Con este nuevo requisito aparece una nueva restricción:
    \[
    x_2 \leq 10 \Rightarrow x_2 + h_4 \geq 0
    \]

    Podemos introducir la nueva restricción en la tabla correspondiente a la solución original:

    \begin{center}
    \begin{tabular}{c|c|cccccccc|}
     & $z$ & $x_1$ & $x_2$ & $x_3$ & $h_1$ & $h_2$ & $h_3$ & $h_4$\\ 
          & $-380$ & $0$ & $0$ & $0$ & $-1$ & $-1$ & $0$ & $0$ \\ 
    \hline
    $x_1$ & 40  & $1$ & $0$ & $0$ & $0$ & $1$ & $1$ & $0$\\ 
    $x_3$ & 20  & $0$ & $0$ & $1$ & $-1$ & $3$ & $-2$ & $0$\\ 
    $x_2$ & 20  & $0$ & $1$ & $0$ & $1$ & $-3$ & $1$ & $0$\\
    \hline
          & 10  & $0$ & $1$ & $0$ & $0$ & $0$ & $0$ & $1$\\
    \hline
    $h_4$ & -10  & $0$ & $0$ & $0$ & $-1$ & $3$ & $-1$ & $1$\\
    \end{tabular}
    \end{center}

    La solución cumple el criterio de optimalidad, pero no es factible, y podemos aplicar Lemke:

    \begin{center}
    \begin{tabular}{c|c|cccccccc|}
     & $z$ & $x_1$ & $x_2$ & $x_3$ & $h_1$ & $h_2$ & $h_3$ & $h_4$\\ 
          & $-380$ & $0$ & $0$ & $0$ & $-1$ & $-1$ & $0$ & $0$ \\ 
    \hline
    $x_1$ & 40  & $1$ & $0$ & $0$ & $0$ & $1$ & $1$ & $0$\\ 
    $x_3$ & 20  & $0$ & $0$ & $1$ & $-1$ & $3$ & $-2$ & $0$\\ 
    $x_2$ & 20  & $0$ & $1$ & $0$ & $1$ & $-3$ & $1$ & $0$\\
    $h_4$ & -10  & $0$ & $0$ & $0$ & $-1$ & $3$ & $-1$ & $1$\\
    \hline
          & $-380$ & $0$ & $0$ & $0$ & $-1$ & $-1$ & $0$ & $0$ \\ 
    \hline
    $x_1$ & 30  &  1 & 0 & 0 &   &  & 0   & \\ 
    $x_3$ & 40  &  0 & 0 & 1 &   &  & 0   & \\  
    $x_2$ & 10  &  0 & 1 & 0 &   &  & 0   & \\ 
    $h_3$ & 10  &  0 & 0 & 0 & 1 & -3 & 1 & -1\\ 
    \hline
    \end{tabular}
    \end{center}

    El nuevo plan de cultivo óptimo pasa por:
    \begin{itemize}
        \item Producir 30 hectáreas de trigo, 40 hectáreas de maíz y 10 hectáreas de girasol.
        \item Se haría uso de toda el agua.
        \item Se emplearía toda la superficie.
        \item Se cumpliría con el requisito ecológico superando en 10 hectáreas el mínimo exigido.
    \end{itemize}
    Esto permite generar una biomasa total de 380 toneladas, igual al plan de cultivo original.


    
\end{enumerate}